\begin{lemma}[7.3.2]
\label{IV.2.s7.3.2-fr}
Sean $X$, $Y$ dos preesquemas localmente noetherianos y sea
$f:X\to Y$ un morfismo. Las propiedades siguientes son equivalentes:
\begin{enumerate}
\item[\emph{a)}] $f$ es un $\boldsymbol P$-morfismo.
\item[\emph{b)}] Para todo $x\in X$, si se pone $y=f(x)$, entonces
$\mathscr O_x$ es un $\mathscr O_y$-módulo plano y se cumple la
propiedad
$\boldsymbol Q(\mathscr O_x\otimes_{\mathscr O_y}\boldsymbol{k}(y),\boldsymbol{k}(y))$.
\item[\emph{c)}] Para todo $x\in X$, si se pone $y=f(x)$, el
morfismo
$\operatorname{Spec}(\mathscr O_x)\to\operatorname{Spec}(\mathscr O_y)$
correspondiente al homomorfismo $\mathscr O_y\to\mathscr O_x$ es un
$\boldsymbol P$-morfismo.
\item[\emph{c$'$)}] Para todo punto cerrado $x\in X$, el morfismo
$\operatorname{Spec}(\mathscr O_x)\to\operatorname{Spec}(\mathscr O_y)$
es un $\boldsymbol P$-morfismo.
\end{enumerate}
\end{lemma}

\begin{proof}
La equivalencia de \emph{a)} y \emph{b)} se deduce de las definiciones; lo mismo
ocurre con la equivalencia de \emph{b)} y \emph{c)}, teniendo en cuenta (\textbf{I}, 2.4.2);
por último, la equivalencia de \emph{b)} y \emph{c$'$)} se deduce de que, para todo
$x\in X$, $\overline{\{x\}}$ contiene un punto cerrado (5.1.11).
\end{proof}
